\section{Life, Movement, and Evolution}

Selves move, live, grow, die, and evolve. Once a self is understood as a pattern rather than its material, all of these become things that happen to patterns, and each gets a clean account. Three questions sharpen this: what exactly is the line between alive and not, how a pattern moves if its cells stay put, and what evolution means when there are no genes at the bottom.

\subsection{How a pattern moves when nothing moves}

Recall the glider. The vibes do not travel. They have no location to travel in. A self moves by re-forming its pattern on neighboring vibes, beat after beat, so the form travels while the vibes stay put, exactly like a wave crossing water without the water going anywhere. Your position is just where your pattern currently sits, and it changes as the pattern propagates.

Three things follow. There is a speed limit: a pattern can re-form no faster than the rule carries influence to the next vibes, which is the light cone, so nothing you are made of can move faster than light. There is inertia: a pattern propagating one way tends to keep doing so, because re-forming in the same direction is the path of least resistance. And your mind goes where your body goes, because your mind is the high-level pattern of your body-region, and it travels with that region. The puzzle of how a moving thing stays the same thing dissolves: it was never the same material, it was always the same pattern.

\subsection{The line between alive and not}

A rock is a stable pattern too, and it is not alive, so stability cannot be the whole story. What is the line? The theory says a living self is a pattern that actively maintains itself, and it names three marks. It keeps a boundary, screening an inside from an outside. It repairs disturbances, healing back to itself when knocked. And it stays open and far from rest, taking in and putting out to hold its form against decay, a flame rather than an ash.

This gives a real criterion for the borderline cases. A flame keeps a boundary and stays far from rest but does not really repair itself in an organized way, so it is lifelike but not alive. A virus keeps a boundary and a pattern but does not maintain itself, it hijacks a cell to do that, so it sits on the border. A crystal is stable but passive, not far from rest, so it is not alive. Life is not a yes-or-no essence but a regime, the boundary-keeping, self-repairing, far-from-equilibrium regime of self, with inert matter at the passive end of the very same spectrum.

\subsection{Growth and death}

A self grows two ways. It grows outward by recruiting neighboring vibes into its pattern, getting bigger. And it grows upward by adding nested sub-selves, climbing the tower, getting more complex. A self dies when its integration drops below the threshold that lets it reconstitute: repair can no longer keep up, the boundary fails, and the pattern disperses. Crucially, the vibes are not destroyed. Only the pattern ends, and the material returns to the field. This is also why, at least in principle, a pattern could reconstitute on fresh material, or re-form from a partial seed, which is the structural picture behind pattern persistence, the idea that what you are is a form that could in principle outlast or recur beyond any particular batch of stuff. The walkthrough flags that as a structural possibility, not a promise.

\subsection{What evolution means here}

Evolution of what, if there are no genes at the bottom? Evolution of patterns. Patterns vary. A pattern that is better at maintaining itself, and better at copying itself onto fresh vibes, lasts longer and spreads, while fragile or sterile patterns vanish. That is variation and selection, the Darwinian ratchet, running on patterns rather than on molecules. Genes are one powerful way for a pattern to copy itself faithfully, but the underlying logic, persist and replicate or disappear, is more general, and it runs on the crystal directly. So in an endlessly growing substrate, the most persistent self-replicating patterns are selected over cosmic time, climbing the tower, and life and mind are simply among the most persistent patterns there are. The qualitative picture is sound. A fully worked model of metabolism and of competing selves is honest future work, not a finished result.

\textbf{Takeaway:} a self moves by re-forming its pattern on new vibes while the vibes stay put, giving a speed limit and inertia for free. Life is not an essence but a regime, the boundary-keeping, self-repairing, far-from-rest kind of self, with viruses and flames sitting believably on the border. Selves grow outward and upward, die when integration fails while the vibes return to the field, and evolve by variation and selection of patterns, the Darwinian logic running directly on the crystal.
